\noindent\textbf{Speculative Decoding.}
Speculative decoding accelerates autoregressive LLM inference by using a draft model to propose candidate tokens and a target model to verify them in parallel. Early approaches use a smaller language model as the drafter, which generates candidate tokens autoregressively before verification by the target model \citep{leviathan2023fast, chen2023accelerating}. Subsequent methods improve this basic draft-then-verify pipeline through tree-based verification, better serving systems, and more efficient draft model designs \citep{miao2023specinfer, cai2024medusa, li2024eagle}. These works establish the general framework of speculative decoding, where the final speedup depends on both the acceptance length and the cost of generating draft tokens.

\noindent\textbf{Autoregressive and Efficient Drafting.}
A representative line of speculative decoding methods improves draft quality through autoregressive drafting. The EAGLE series generates draft tokens sequentially, allowing each token to depend on previous draft tokens and better match the target model's autoregressive distribution \citep{li2024eagle, li2024eagle2, li2025eagle3}. Other methods reduce drafting overhead from different angles: Medusa uses lightweight parallel decoding heads \citep{cai2024medusa}, Hydra introduces sequentially-dependent heads to inject causal information into head-based drafting \citep{ankner2024hydra}, and FR-Spec and SpecVocab reduce full-vocabulary projection costs through static or dynamic vocabulary selection \citep{zhao2025fr,williams2026speculative}. 

\noindent\textbf{Parallel and Non-Autoregressive Drafting.}
Speculative Diffusion Decoding first explores discrete diffusion models as parallel drafters for speculative decoding \citep{christopher2025speculative}. DiffuSpec further uses pretrained diffusion language models as training-free drafters \citep{li2025diffuspec}. PARD adapts autoregressive models into parallel draft models, allowing multiple future tokens to be predicted in a single draft forward pass \citep{an2026pard}.  More recently, DART predicts token distributions for multiple future positions in parallel and uses  tree pruning to construct draft candidates \citep{liu2026dart}. DFlash adopts a block-diffusion drafter that produces an entire draft block in a single forward pass, avoiding repeated calls to both the draft model and the full LM head \citep{chen2026dflash}. These methods substantially reduce drafting overhead, but fully parallel drafting weakens intra-block causal dependencies, making it harder to match the target model's autoregressive distribution. In contrast, Domino keeps the main draft computation parallel while reintroducing causal information through a lightweight correction branch.